% =============================================
% CHƯƠNG 5 — TRIỂN KHAI VÀ KIỂM THỬ (9 trang)
% Status: OUTLINE — chờ draft từ agent
% =============================================

\chapter{Triển khai và kiểm thử}
\label{ch:trien-khai}

% [INTRO CHƯƠNG]
% Chương này mô tả chi tiết quá trình triển khai từ thiết kế ở chương 4 thành
% sản phẩm thực tế, bao gồm các module Mini App, dashboard quản trị, Cloud Functions,
% kết quả kiểm thử, đánh giá bảo mật và hiệu năng, cũng như quy trình phát hành.

% ─────────────────────────────────────────────
\section{Triển khai Mini App}
\label{sec:trien-khai-miniapp}
% Ước lượng: 1.5 trang
% Nguồn: code src/ + screenshots (CẦN CHỤP qua Playwright)
% Nội dung:
%   §5.1.1 Cấu trúc thư mục (98 files, ~15.9K dòng)
%   §5.1.2 Các trang chính:
%     - splash, welcome, login (với 3 bước MSSV → email → OTP)
%     - home (dashboard, calendar, menu)
%     - student/* (6 trang)
%     - teacher/* (7 trang)
%     - AIChatPage
%   §5.1.3 Module dùng chung: 10 hooks, 9 UI components
%   §5.1.4 Tích hợp Zalo SDK (authorize, getUserInfo, scanQRCode, ...)
%   Hình 5.1: Cấu trúc thư mục mini app
%   Hình 5.2: Screenshot 4 màn hình chính (splash, login, home, attendance)

% ─────────────────────────────────────────────
\section{Triển khai Admin Dashboard}
\label{sec:trien-khai-admin}
% Ước lượng: 1 trang
% Nguồn: báo cáo 23-28/03 + admin/src/pages
% Nội dung:
%   §5.2.1 Tech stack: Vite + React + Ant Design + Recharts + xlsx
%   §5.2.2 12 trang chính: Dashboard, Users, Classes, ClassDetail, Attendance,
%     AbsenceRequests, FraudReports, PresentPage, ...
%   §5.2.3 Tính năng đặc biệt: import Excel SV/Lớp, export Excel báo cáo, charts
%   §5.2.4 Phân quyền: Firebase Auth email/password (tách biệt Zalo OAuth)
%   §5.2.5 Deploy: Firebase Hosting `inhust-admin.web.app`
%   Hình 5.3: Screenshot Dashboard tổng quan + một trang chi tiết

% ─────────────────────────────────────────────
\section{Triển khai Cloud Functions}
\label{sec:trien-khai-functions}
% Ước lượng: 1 trang
% Nguồn: functions/src/index.ts (17 hàm)
% Nội dung:
%   §5.3.1 Danh sách 17 hàm theo module:
%     - session: createSession, endSession
%     - attendance: submitAttendance, manualAttendance
%     - trust: calculateTrustScore
%     - fraud: detectFraud, generateFraudReport
%     - face: registerFace, verifyFace
%     - auth: createCustomToken, oauthCallback
%     - admin: dashboardStats, bulkAssignStudents, reviewAbsenceRequest, createAccount
%     - ai: aiChat (Groq proxy)
%   §5.3.2 Pattern chung: httpsCallable, requireAuth middleware
%   §5.3.3 Hạn chế Spark plan: không deploy được (chỉ test emulator)
%   Bảng 5.1: 17 Cloud Functions với mô tả ngắn

% ─────────────────────────────────────────────
\section{Kiểm thử}
\label{sec:kiem-thu}
% Ước lượng: 2 trang
% Nguồn: docs/test-scenarios.md (29 TC) + báo cáo các tuần
% Nội dung:
%   §5.4.1 Chiến lược kiểm thử:
%     - Unit test (TypeScript): Jest cho utility functions
%     - Integration test: emulator Firebase
%     - E2E test: Playwright (kế hoạch, chưa triển khai đầy đủ)
%     - Manual test: 29 test case theo kịch bản
%   §5.4.2 Kết quả unit test
%   §5.4.3 Kết quả manual test (29 TC):
%     Bảng 5.2 — TC | Mô tả ngắn | Kỳ vọng | Kết quả | Pass/Fail
%     (lấy từ test-scenarios.md, tổng hợp gọn 15-20 hàng quan trọng)
%   §5.4.4 Bug đã phát hiện và sửa (chọn 3-5 bug có giá trị):
%     - Race condition useEffect dependency (báo cáo 11-16/05)
%     - Wall-clock sync between devices (báo cáo 11-16/05)
%     - GPS pre-warm timing (báo cáo 04-09/05)
%     - Firebase initializeApp twice (báo cáo 18-23/05)
%     - Privacy Policy 404 (báo cáo 18-23/05)

% ─────────────────────────────────────────────
\section{Đánh giá bảo mật}
\label{sec:bao-mat}
% Ước lượng: 1.5 trang
% Nguồn: docs/SECURITY-REPORT.md (đã có 7.6/10)
% Nội dung:
%   §5.5.1 Phương pháp đánh giá: rubric 10 lĩnh vực
%   §5.5.2 Bảng 5.3 — Điểm 10 lĩnh vực bảo mật:
%     1. Authentication, 2. Authorization, 3. Data validation,
%     4. Encryption, 5. Session management, 6. API security,
%     7. Privacy, 8. Audit logging, 9. Rate limiting, 10. Compliance
%   §5.5.3 Các lỗ hổng đã fix (lấy từ SECURITY-REPORT "DA FIX"):
%     - Groq API key client-side → Cloud Function proxy
%     - HMAC secret client → server-side
%     - Firestore rules `if true` → role-based
%     - Pairing token guess → 128-bit random + one-shot
%   §5.5.4 Hạn chế còn lại (lấy từ "Chua lam duoc - Dai han"):
%     - Spark plan chưa có rate limiting backend
%     - Face anti-spoofing chưa có liveness
%     - Cần audit log dài hạn

% ─────────────────────────────────────────────
\section{Đánh giá hiệu năng}
\label{sec:hieu-nang}
% Ước lượng: 1 trang
% VIẾT MỚI — cần đo thực tế
% Nội dung:
%   §5.6.1 Phương pháp đo: Lighthouse, React DevTools, Firestore usage console
%   §5.6.2 Bảng 5.4 — Metrics đo được:
%     - Cold start Mini App: __ms
%     - Render splash → home: __ms
%     - QR scan to verification: __s
%     - Firestore reads per session: __ docs
%     - Bundle size: __ KB
%   §5.6.3 Bottleneck đã xử lý:
%     - N+1 query trong TeacherClassDetail (báo cáo 04-09/05)
%     - GPS pre-warm để giảm wait time
%   §5.6.4 Hướng tối ưu tiếp theo (đề xuất cho phần kết luận)

% ─────────────────────────────────────────────
\section{Quy trình phát hành và tuân thủ}
\label{sec:phat-hanh}
% Ước lượng: 1 trang
% Nguồn: báo cáo 18-23/05 + legal/*.html
% Nội dung:
%   §5.7.1 Đăng ký Mini App trên Zalo Platform
%   §5.7.2 Tuân thủ Zalo Censorship Policy:
%     - Privacy Policy 11 mục
%     - Terms of Service
%     - Consent dialog AI Chat
%     - Quyền xoá dữ liệu khuôn mặt
%   §5.7.3 Quyền API được cấp:
%     - OA Send Message: 13 quyền duyệt
%     - ZNS API: 3 quyền duyệt
%   §5.7.4 Hạn chế khi triển khai chính thức:
%     - Cần OA verified để dùng ZNS production
%     - Cần Quyết định thành lập trường để vào danh mục Giáo dục
%   §5.7.5 Phương án cho phạm vi đồ án: danh mục "Công cụ", demo qua sandbox

% ─────────────────────────────────────────────
% \section*{Kết chương}
% Đã triển khai đầy đủ thiết kế chương 4 và đánh giá theo 4 chiều: chức năng,
% bảo mật, hiệu năng, tuân thủ. Chương cuối tổng kết kết quả và đề xuất phát triển.
